\chapter{Conclusion}
In this master thesis, we have presented a introduction of Martin Hairer’s theory of regularity structures, and then focusing on the modern diagram-free approach to the stochastic quantization of the $\phi^4_3$ model. 
By substituting the intricate combinatorial machinery of decorated trees with a parsimonious index set of multi-indices, this framework streamlines the proof of convergence. The key insight lies in treating the directional Malliavin derivative as a tangent vector to the solution manifold, which allows for a controlled representation as a modeled distribution, with a gain in regularity with respect to the components of the model. 
Combined with the spectral gap inequality, this geometric formulation reduces the stochastic challenge of bounding the divergent model components to a robust, pathwise analytic control of Malliavin increments where counterterms completely vanish. 
Throughout this work, we have detailed the core analytical steps of reconstruction and integration. 
Moving forward, a natural next step would be to better undersand the bridges between this approach and the tree-based approach. Notably, having better interpretation of the coefficients $\dG$.
